%% kapa: exam=1 round=22 year=2011 period=2 form=A subject=appraisal_law subject_ko=부동산관계법규 seq=18 number=58
%% source: Q-Net 감정평가사 기출문제 (cst00309) artlSeq=5036790 file=extracted/1008/5036790/2011년 제22회 감정평가사 1차시험 기출문제/2교시.hwp
%% conversion: kordoc 4.24.51 -> tools/gen1.py (auto); review=proofread-gpt-astra; flags=none
\begin{question}{58}
부동산 가격공시 및 감정평가에 관한 법령상 감정평가사에 대한 징계에 관한 설명으로 옳은 것은?
\begin{choices}
  \item 감정평가협회는 감정평가사징계위원회에 감정평가사에 대한 징계의결의 요구를 할 수 없다.
  \item 감정평가업자가 중과실로 잘못된 평가를 한 경우에도 고의가 없다면 징계의 대상이 되지 않는다.
  \item 감정평가사징계위원회의 의사정족수는 재적위원 과반수이고 의결정족수는 출석위원 3분의 2 이상이다.
  \item 감정평가사가 업무상 중대한 의무를 위반한 경우 감정평가사징계위원회의 의결에 따라 자격취소의 징계를 할 수 있다.
  \item 감정평가사에 대한 징계의결의 요구는 위반사유가 발생한 날부터 3년이 지난 때에는 할 수 없다.
\end{choices}
\end{question}
\answer{1}
